% !TEX root = ../presentation.tex
% =============================================================
%  Symbols slide
%  Part of ac-transmission; \input from ac-transmission.tex.
%  Two columns, same layout as the Abbreviations frame.
% =============================================================

\begin{frame}{Symbols}
  \renewcommand{\arraystretch}{1.25}
  \begin{columns}[T,onlytextwidth]
    \begin{column}{0.48\textwidth}
      {\small
      \begin{tabularx}{\linewidth}{@{}p{0.9cm}>{\raggedright\arraybackslash}X@{}}
        \textbf{Sym.} & \textbf{Description} \\
        \hline
        $P$      & Active power \\
        $V_{LL}$ & Line-to-line (phase-phase) voltage \\
        $V_{LN}$ & Line-to-neutral (phase-ground) voltage \\
        $\delta$ & Transmission angle between the terminal voltages \\
        $X$      & Series reactance of the path between the terminals \\
        $L'$     & Series inductance per unit length \\
      \end{tabularx}}
    \end{column}
    \begin{column}{0.48\textwidth}
      {\small
      \begin{tabularx}{\linewidth}{@{}p{0.9cm}>{\raggedright\arraybackslash}X@{}}
        \textbf{Sym.} & \textbf{Description} \\
        \hline
        $D_m$    & Geometric mean distance between the phases \\
        $r_b$    & Equivalent radius of the conductor bundle \\
        $Z_w$    & Wave (surge) impedance of the line \\
        $\beta$  & Phase constant, $\beta\cdot l$ the electrical length \\
        $K_s$    & Synchronising power coefficient, $\partial P/\partial\delta$ \\
        $\underline{\gamma}$ & Propagation constant, $\alpha+j\cdot\beta$ \\
        $\alpha$ & Attenuation constant \\
        $H$      & Inertia constant of a synchronous machine \\
      \end{tabularx}}
    \end{column}
  \end{columns}
\end{frame}
